% ============================================================
\section{Problem Statement}
\label{sec:problem-statement}
% ============================================================

This section formalises the problem addressed by the present work. We first
explain why the inventory of advertising formats inherited from the visual web
fails to transfer to conversational interfaces, then enumerate the technical
sub-problems that any LLM-native advertising system must solve, then state the
ethical constraints that bound the admissible solution space, and finally
present a precise formal definition of the task.

\subsection{Limitations of Traditional Ad Formats in Conversational AI}
\label{subsec:traditional-ad-limitations}

Three families of advertising formats dominate the contemporary web, and each
fails to transfer to a conversational setting for a structural reason rather
than an incidental one.

\emph{Display advertising}---banners, sidebars, and in-feed units---presupposes
a two-dimensional visual canvas in which a rectangular region can be reserved
for sponsored content without disturbing the primary content stream. A chat
interface offers no such region. The output of an LLM is a single linear
sequence of tokens rendered as flowing text, and any attempt to reintroduce a
sidebar or banner reasserts the very visual paradigm that conversational
systems were designed to replace.

\emph{Interruption-based formats}---pre-roll video, interstitials, and modal
overlays---monetise attention by suspending the user's task for a fixed
duration. In a conversational exchange, an interruption between the user's
question and the assistant's answer breaks the implicit contract that the
system responds to what was asked. The cost of such an interruption is not
merely an aesthetic complaint; it directly degrades the property---fluid
turn-taking---that makes conversational interfaces valuable in the first place.

\emph{Search advertising} appears at first glance to be the closest analogue,
since both search and chat are text-driven and intent-rich. The analogy is
nevertheless misleading. A search query is a short, atomic keyword expression
matched against a bid space by exact and approximate string operators. A
conversational query is multi-turn: the relevant intent is distributed across
several user messages, refined by clarifying questions, and frequently never
expressed as a single noun phrase. A keyword auction operating on the most
recent user message alone discards the bulk of the signal that the conversation
contains.

The conclusion is that advertising in conversational AI is not a porting
problem. It requires a format that is itself conversational---one in which
sponsored content is integrated into the assistant's reply as a natural
recommendation rather than appended as a foreign object.

\subsection{Technical Challenges}
\label{subsec:technical-challenges}

Designing such a format raises four technical sub-problems that must be solved
jointly.

\subsubsection{Real-Time Intent Extraction from Multi-Turn Conversations}
\label{subsubsec:intent-extraction}

Let $H_t = (m_1, m_2, \ldots, m_{t-1})$ denote the conversation history up to
turn $t$ and $m_t$ the current user message. The system must produce, in real
time, a structured representation $C_t = \phi(H_t, m_t)$ that captures the
user's latent commercial intent, the topical focus of the conversation, the
strength of any purchase signal, and the user's receptivity to advertising at
the present moment. The mapping $\phi$ cannot be reduced to keyword extraction:
purchase intent often emerges only when a sequence of exploratory questions is
considered jointly, and receptivity depends on conversational register
(casual chit-chat, technical research, emotional disclosure) rather than on any
single lexical cue. The extraction must additionally complete within a latency
budget compatible with interactive use, ruling out heavyweight inference.

\subsubsection{Semantic Ad Matching at Low Latency}
\label{subsubsec:semantic-matching}

Given the extracted context $C_t$ and an ad inventory
$\mathcal{A} = \{a_1, \ldots, a_N\}$, the system must select a small set of
candidate advertisements that are semantically aligned with the conversation.
Lexical matching against ad metadata is insufficient: a user discussing
``training for a half marathon'' should surface running shoes even though the
phrase ``running shoes'' never appears in the dialogue. The matching function
must therefore operate in a semantic space, must respect business constraints
(active campaigns, positive remaining budget, frequency caps), and must be
computable within a few hundred milliseconds over inventories that may grow to
tens of thousands of items. A composite scoring rule is required, in which
semantic similarity, advertiser bid, and the inferred purchase signal are
combined into a single ranking criterion whose weights are interpretable and
auditable.

\subsubsection{Natural Ad Integration Without Degrading User Experience}
\label{subsubsec:natural-integration}

Even a perfectly matched advertisement is worthless if it is appended to the
response as a foreign block. The selected ad $a^\star$ must be woven into the
assistant's reply as a recommendation that flows from the surrounding text,
while a mandatory disclosure marker---we adopt the literal token
\texttt{[Sponsored]}---must accompany every commercial mention to satisfy
regulatory requirements. Two failure modes must be excluded by construction:
the assistant must not promote $a^\star$ when no candidate is genuinely
relevant to the user's question, and it must not produce a response in which
the disclosure is omitted, hidden, or rephrased. Both constraints place a
non-trivial burden on the prompt that drives the response-synthesis stage,
since the model must be simultaneously encouraged to recommend and permitted
to refuse.

\subsubsection{Engagement Measurement Beyond Clicks}
\label{subsubsec:engagement}

Click-through rate, the canonical metric of digital advertising, is poorly
suited to conversational interfaces. Clicks are sparse---most users do not
leave the conversation to visit an external page---and a click measures only
the act of departure, not the quality of attention. A more faithful signal is
available in the conversation itself: when a user follows a sponsored mention
with a question about the advertised product, that follow-up is direct
evidence that the advertisement entered the user's working model of the
exchange. Detecting such follow-ups requires a classifier that takes both the
prior advertisement and the next user message as input, and that distinguishes
genuine engagement from coincidental topical overlap. We treat this engagement
event as a first-class metric, complementary to and stronger than the click.

\subsection{Ethical and Regulatory Considerations}
\label{subsec:ethical-considerations}

Conversational advertising operates in a setting of elevated user trust. A
chat assistant is perceived---accurately or not---as a neutral helper, and any
recommendation it issues carries the weight of that perceived neutrality. This
raises three constraints that the system must respect.

\emph{Disclosure.} United States Federal Trade Commission guidance requires
that sponsored content be identified clearly and conspicuously. We adopt the
explicit token \texttt{[Sponsored]} adjacent to every commercial mention, and
treat the absence of this token as a hard failure of the response-synthesis
stage rather than a stylistic preference.

\emph{Sensitivity gating.} Certain conversational contexts---disclosures of
illness, bereavement, financial distress, or acute emotional difficulty---are
incompatible with advertising of any kind. The context-extraction stage must
therefore emit not only a positive receptivity score but also an explicit
\emph{none} value that causes the matching pipeline to be bypassed entirely.
This is treated as an inviolable safeguard: under sensitivity gating, no
advertisement is shown regardless of bid, similarity, or budget.

\emph{Data minimisation.} Session state is retained only for the duration of
an active conversation, with a fixed inactivity expiry. No persistent user
profile is constructed across sessions, and no session content is shared with
advertisers beyond the aggregate impression and engagement counters required
for billing.

These constraints are not soft preferences. They define the admissible region
of the design space and any solution that violates them is rejected
\emph{a priori}, irrespective of its performance on other metrics.

\subsection{Formal Problem Definition}
\label{subsec:formal-problem}

We may now state the problem precisely. Let $\mathcal{M}$ denote the space of
natural-language messages and $\mathcal{A}$ a finite ad inventory. At turn $t$
the system observes a history $H_t \in \mathcal{M}^{t-1}$ and a user message
$m_t \in \mathcal{M}$, and must produce a tuple
\[
    \big(r_t,\; \alpha_t,\; e_t\big)
    \;\in\; \mathcal{M} \times \big(\mathcal{A} \cup \{\varnothing\}\big)
                       \times \mathcal{E},
\]
where $r_t$ is the assistant's reply, $\alpha_t$ is the advertisement embedded
in $r_t$ (or $\varnothing$ if none is shown), and $e_t$ is the set of
engagement events logged for downstream analytics. The tuple is required to
satisfy the following conditions.

\begin{description}[leftmargin=2.2em,style=nextline]
    \item[(C1) Helpfulness.] The reply $r_t$ answers the user's query in $m_t$
        on its own merits, independently of whether an advertisement is
        included.
    \item[(C2) Relevance-conditional injection.] If $\alpha_t \neq \varnothing$,
        then $\alpha_t$ is semantically aligned with the inferred context
        $C_t = \phi(H_t, m_t)$ and lies within the active, budgeted, and
        frequency-cap-eligible subset of $\mathcal{A}$.
    \item[(C3) Mandatory disclosure.] If $\alpha_t \neq \varnothing$, then
        $r_t$ contains the disclosure marker \texttt{[Sponsored]} adjacent to
        the mention of $\alpha_t$.
    \item[(C4) Sensitivity gating.] If $C_t$ indicates a sensitive context,
        then $\alpha_t = \varnothing$ and no element of $\mathcal{A}$ is
        evaluated for inclusion.
    \item[(C5) Latency.] The end-to-end production of $(r_t, \alpha_t, e_t)$
        completes within a fixed wall-clock budget compatible with interactive
        chat.
\end{description}

The objective is to design a system that satisfies (C1)--(C5) on every turn
while maximising a composite utility that rewards both user-perceived
helpfulness and verifiable advertiser engagement. The remainder of this report
develops such a system, presents the architectural decisions that realise each
of (C1)--(C5), and evaluates the resulting platform on a curated set of
representative conversations.